%%%%%%%%%%%%%%%%%%%%%%%%%%%%%%%%%%%%%%%%%%%%%%%%%%%%%%%%%%%%%%%%%%%%%%%%%%%%%%%
% SECTION 3: Backpropagation
%%%%%%%%%%%%%%%%%%%%%%%%%%%%%%%%%%%%%%%%%%%%%%%%%%%%%%%%%%%%%%%%%%%%%%%%%%%%%%%
\section{Backpropagation}

% Feedforward Neural Network Structure
\begin{frame}{Feedforward Neural Network Structure}
    % TODO: Add diagram of feedforward NN with labels:
    % - Input layer (x₁, x₂, ..., xₙ)
    % - Hidden layers with neurons
    % - Weights (wᵢⱼ) on connections
    % - Biases (bⱼ) at each neuron
    % - Output layer
    
    \begin{center}
        \textit{[Feedforward NN diagram with labeled weights $w_{ij}$ and biases $b_j$]}
    \end{center}
    
    \vspace{1em}
    \textbf{Each neuron computes:}
    $$z_j = \sum_i w_{ij} x_i + b_j \quad \text{then} \quad a_j = \sigma(z_j)$$
    
    where $\sigma$ is the activation function
\end{frame}

% The Algorithm
\begin{frame}{Backpropagation: The Algorithm}
    \textbf{The training algorithm that makes representation learning practical}
    
    \vspace{1em}
    \begin{enumerate}
        \item \textbf{Forward pass:} compute prediction $\hat{y}$
        \item \textbf{Compute loss:} $L = \frac{1}{2}(y - \hat{y})^2$
        \item \textbf{Backward pass:} compute gradients via chain rule
        \item \textbf{Update weights:} $w \leftarrow w - \eta \frac{\partial L}{\partial w}$
        \item Repeat
    \end{enumerate}
\end{frame}

% Backprop Math
\begin{frame}{Backpropagation: The Math}
    \textbf{Chain rule propagates error backward:}
    
    \vspace{1em}
    For weight $w_{ij}$ connecting layer $i$ to layer $j$:
    $$\frac{\partial L}{\partial w_{ij}} = \frac{\partial L}{\partial a_j} \cdot \frac{\partial a_j}{\partial z_j} \cdot \frac{\partial z_j}{\partial w_{ij}}$$
    
    \vspace{1em}
    \textbf{Key insight:} Each term is local
    \begin{itemize}
        \item $\frac{\partial z_j}{\partial w_{ij}} = a_i$ (input from previous layer)
        \item $\frac{\partial a_j}{\partial z_j} = \sigma'(z_j)$ (derivative of activation)
        \item $\frac{\partial L}{\partial a_j}$ propagates from layers above
    \end{itemize}
    
    \vspace{0.5em}
    \textit{This is why activation functions must be differentiable}
\end{frame}

% Stochastic Gradient Descent
\begin{frame}{Stochastic Gradient Descent (SGD)}
    \textbf{Why ``stochastic''?}
    
    \vspace{1em}
    \textbf{Batch gradient descent:}
    $$w \leftarrow w - \eta \frac{1}{N} \sum_{i=1}^{N} \nabla_w L_i$$
    Compute gradient over \textit{all} $N$ training examples — slow!
    
    \vspace{1em}
    \textbf{Stochastic gradient descent:}
    $$w \leftarrow w - \eta \nabla_w L_i \quad \text{(single random example)}$$
    
    \vspace{0.5em}
    \textbf{Mini-batch SGD:} Use small random subset (e.g., 32 examples)
    \begin{itemize}
        \item Noisy estimate of true gradient
        \item But much faster iteration
        \item Noise can help escape local minima
    \end{itemize}
\end{frame}

% Activation Functions
\begin{frame}{Activation Functions}
    \textbf{Why do we need them?}
    \begin{itemize}
        \item Without non-linearity, stacked layers collapse to single linear transform
        \item Non-linear $\sigma$ enables learning complex functions
    \end{itemize}
    
    \vspace{1em}
    \textbf{Sigmoid / Tanh:}
    \begin{itemize}
        \item Smooth, differentiable
        \item Problem: cause \textbf{vanishing gradients} — $\sigma'(z) \to 0$ for large $|z|$
    \end{itemize}
    
    \vspace{0.5em}
    \textbf{ReLU} $(\max(0, x))$:
    \begin{itemize}
        \item Simple, fast to compute
        \item Gradient is 1 for $x > 0$ — no vanishing
        \item The paper's recommended default
    \end{itemize}
    
    % TODO: Add activation function plots (sigmoid, tanh, ReLU)
\end{frame}

% Code Example Placeholder
\begin{frame}[fragile]{Backpropagation: Code Example}
    % TODO: Add simple neural network code example
    % - Define a 2-layer network
    % - Forward pass
    % - Loss computation  
    % - Backward pass (manual gradients)
    % - Weight update
    
    \begin{center}
        \textit{[Code: simple 2-layer network with ReLU, forward/backward pass]}
    \end{center}
\end{frame}

% The 1990s Objection
\begin{frame}{The 1990s Objection}
    \textbf{The prevailing concern:}
    \begin{itemize}
        \item Backpropagation would get trapped in poor \textbf{local minima}
        \item Points where gradient is zero but solution is far from optimal
        \item Widely cited as proof deep networks could never work
    \end{itemize}
    
    \vspace{1em}
    \textbf{The paper's rebuttal:}
    \begin{quote}
        \small``The landscape is packed with a combinatorially large number of saddle points where the gradient is zero, and the surface curves up in most dimensions and curves down in the remainder.''
    \end{quote}
    
    \vspace{0.5em}
    \textbf{Saddle points are not local minima} — there is always a descent direction
\end{frame}

% Feedforward NN Weights Visualization
\begin{frame}{What Do Feedforward Networks Learn?}
    % TODO: Add visualization of learned weights from a simple feedforward NN
    % trained on images - show that weights look somewhat random/unstructured
    % compared to CNN filters
    
    \begin{center}
        \textit{[Visualization: weight matrices from feedforward NN on images]}
    \end{center}
    
    \vspace{1em}
    \textbf{Observation:} Weights appear relatively unstructured
    \begin{itemize}
        \item No clear feature detectors
        \item Contrast with CNNs which learn interpretable filters
    \end{itemize}
    
    \vspace{0.5em}
    \textit{This motivates architectures with built-in structure...}
\end{frame}
